Now that we have a unifying framework in which many different libraries are represented, we can approach the problem of integrating these libraries with each other, e.g. to translate knowledge between them. \medskip

Naturally, these libraries use different foundations; which is reflected in \mmt as the contained theories having different meta theories. Consequently, as a first step one might want to integrate the foundations themselves, i.e. by implementing systematic translations between them. 

An obvious challenge arises here, namely that the foundations used by formal systems are usually mutually incompatible. This is obvious for systems using fundamentally different logics, such as set theories vs. type theories, but even if the \emph{platonic} foundation used by two systems is the same (e.g. higher-order logic), there are subtle differences in the specific implementations -- and hence ontologies -- of the foundations, that make implementing such translations less straight-forward than one might at first suspect. For example, the systems might offer different additional features (e.g. records, subtyping) where there is no universally valid elimination procedure. Additionally, systems usually offer distinct convenience operations, automations, module systems and conventions, that are suggestive of different approaches and best practices for formalizing content, leading to very different library developments for the same mathematical concept even if the foundation used by the systems is the same.

As a result, even when a translation between foundations is straight-forward (which it usually is not), the result of translating content between libraries \emph{on the basis of} a translation between the foundations can yield awkward reformalizations of known content completely disconnected from the content already available in the target library.

Consequently, any generic methodology for integrating libraries needs to go beyond -- or even entirely bypass -- mere foundations and take distinct library developments into account. This part of this thesis presents our approach to library integration on that premise.\medskip

First, we introduce the notion of \emph{alignments} in \autoref{sec:crosslib:alignments}. This is a binary relation between symbols expressing that the two symbols (possibly from different libraries) denote the same abstract mathematical concept. In the simplest case, we can then translate symbol occurences in expressions across libraries by merely substituting along alignments. We then classify various more complex types of alignments (e.g. ``up to argument order'' or ``up to additional arguments'') that occur in practice and develop a convenient way to specify them.

\autoref{sec:crosslib:translate} generalizes this approach to translating expressions; using alignments wherever possible, but additionally using theory morphisms (within or across libraries) and -- as is crucial for translating across different foundations -- programmatic translations, which can be supplied for more complex situations. 

To avoid having to provide ad-hoc traslation mechanisms between all pairs of libraries, we instead try to use translations between any library and a fixed set of \emph{interface theories}, in line with the Math-in-the-Middle approach (see \autoref{sec:lf:mitm}). In a library of interface theories, we can implement all possible developments and definitions of the same mathematical concept and connect them via theory morphisms, lifting the more complicated translation problems to a realm governed by a single, flexible foundation using the logical frameworks developed in \autoref{ch:lf}. Not surprisingly, this methodology is used to implement the functionalities aimed at by the \ODK project (see \autoref{sec:intro:prelim:odk}) -- \autoref{sec:crosslib:translate:ex} shows a real-world use case from the \ODK community that is covered by our approach.

In \autoref{sec:crosslib:vf}, we describe how to automatically find theory morphisms; both within and across different libraries. These can be used to identify overlap between libraries and consequently suggest new alignments and translations between the theories; ideally resulting in a positive feedback loop of ever increasing connections between libraries.

Lastly, \autoref{sec:crosslib:intersect} develops a method for using theory morphisms to refactor library content to increase modularity within a given library.